\documentclass[11pt,letterpaper]{article}
\usepackage{../assets/scientific_report}
\usepackage{float}
\usepackage{listings}
\usepackage{xcolor}
\usepackage{graphicx}
\usepackage{geometry}
\geometry{headheight=23pt, top=2cm, bottom=2cm, left=2.5cm, right=2.5cm}
\setlength{\parskip}{0.5em}
\setlength{\parindent}{0pt}
\lstset{
  language=C,
  basicstyle=\ttfamily\small,
  keywordstyle=\color{blue}\bfseries,
  commentstyle=\color{green!60!black}\itshape,
  stringstyle=\color{red!80!black},
  numbers=left,
  numberstyle=\tiny\color{gray},
  numbersep=8pt,
  frame=single,
  breaklines=true,
  captionpos=b,
  tabsize=4,
  showstringspaces=false
}
\begin{document}
\pagestyle{plain}
\begin{center}
{\LARGE\bfseries Regioes Criticas Nomeadas e Locks Explicitos}\\[0.4em]
{\normalsize Eugenio Vitor Lopes dos Santos}\\[0.3em]
{\small DCA3703, Programacao Paralela, UFRN}\\[0.3em]
\end{center}
\section{Enunciado}
Escrever um programa que cria tarefas para realizar N insercoes em duas listas encadeadas. Cada tarefa escolhe aleatoriamente em qual lista inserir. Garantir integridade evitando condicao de corrida e, sempre que possivel, usar regioes criticas nomeadas para que a insercao em uma lista nao bloqueie a outra. Generalizar para numero de listas definido pelo usuario. Explicar por que regioes criticas nomeadas nao bastam e por que locks explicitos se tornam necessarios.
\section{Fundamentacao}
Inserir no inicio custa dois stores: \texttt{novo->prox = head; head = novo}. Duas threads lendo o mesmo \texttt{head} perdem uma insercao (\textit{lost update}).

\texttt{critical} sem nome trava tudo num lock global unico, inclusive listas independentes. \texttt{critical(name)} separa um lock por nome: mesmo nome espera junto, nome diferente anda junto.

O nome sai resolvido em compilacao e nao aceita variavel como \texttt{critical(locks[i])}. Para N definido em execucao foi alocado \texttt{omp\_lock\_t locks[N]}: \texttt{omp\_init\_lock} antes do paralelo, \texttt{omp\_set\_lock}/\texttt{omp\_unset\_lock} por insercao, \texttt{omp\_destroy\_lock} depois. Foram criadas \texttt{tasks} com \texttt{parallel+single+task} e foi sorteado com \texttt{rand\_r(\&seed)} privada por tarefa; \texttt{rand()} divide estado global e corre com concorrencia.
\section{Implementacao}
Foram escritos dois arquivos no padrao v1/v2 da tarefa08 e compilados com GCC 13.3, \texttt{-std=c11 -O2 -Wall -fopenmp -lm}. Foi medido com \texttt{agora()}, sobre \texttt{clock\_gettime(CLOCK\_MONOTONIC)}: resolucao de nanossegundos, imune a ajustes do relogio.

Cada cronometro cobre so a regiao paralela de insercoes (\texttt{parallel+single+tasks}). Alocacao, \texttt{omp\_init\_lock} e \texttt{omp\_destroy\_lock} ficam de fora. Foi fixado o escopo com \texttt{default(none)}. v1 guarda duas cabecas com \texttt{critical(lista0)}/\texttt{critical(lista1)}; v2 guarda \texttt{heads[nlistas]} com \texttt{locks[nlistas]}. Foram contados os nos no fim e conferido \texttt{total==N}.
\begin{lstlisting}[language=bash,caption={Compilacao e execucao}]
gcc -std=c11 -O2 -Wall -fopenmp v1_critical_nomeado.c -o v1_critical -lm
gcc -std=c11 -O2 -Wall -fopenmp v2_locks.c -o v2_locks -lm
./v1_critical 50000
./v2_locks 50000 8
OMP_NUM_THREADS=4 ./v2_locks 100000 8
\end{lstlisting}
\section{Resultados}
Foram executados testes num i7-14700HX sob WSL2 (14 nucleos fisicos, 28 threads logicas). Foi varrido 1 a 28 threads com N = 50000 e 8 listas na v2. Cada ponto vale a mediana de 3 repeticoes. CSV completo em anexo.
\begin{lstlisting}[language=bash,caption={Saidas reais de integridade}]
[v1 critical nomeado] N=50000 c0=24998 c1=25002 total=50000 OK tempo=0.3221 s
[v2 locks] N=50000 nlistas=8 total=50000 OK tempo=0.3807 s
[v2 locks] N=50000 nlistas=2 total=50000 OK tempo=0.3531 s
\end{lstlisting}
\begin{figure}[H]
\centering
\includegraphics[width=0.9\textwidth]{tempo_execucao_plot.png}
\caption{Tempo contra threads em escala linear.}
\end{figure}
\begin{figure}[H]
\centering
\includegraphics[width=0.6\textwidth]{comparacao_versoes.png}
\caption{Comparacao v1 x v2.}
\end{figure}
\section{Discussao}
\subsection{Sem bloqueio cruzado na v1}
\texttt{critical} sem nome trava tudo num lock global: a lista 0 bloqueia a 1 sem dividir dado. Foram separados \texttt{critical(lista0)} e \texttt{critical(lista1)} em dois \texttt{locks} independentes. Mesma lista espera junto; listas diferentes andam juntas. Sobra so a espera inevitavel: duas tarefas na mesma lista.
\subsection{Por que o nomeado nao generaliza}
O nome do critical e literal de compilacao e nao aceita variavel:
\begin{lstlisting}[language=C,caption={Nao indexavel: nome precisa ser literal}]
// invalido: nao compila / nao indexa por i
#pragma omp critical(locks[i])
  inserir(&heads[i], v);
\end{lstlisting}
1000 listas pediriam 1000 nomes escritos a mao, e N so chega em execucao via \texttt{argv[2]}. Um nome unico para tudo volta a serializar o que e independente. Foi alocado \texttt{malloc + omp\_init\_lock} antes do paralelo, foi travado \texttt{omp\_set\_lock}/\texttt{omp\_unset\_lock} por \texttt{qual} a cada insercao e liberado com \texttt{omp\_destroy\_lock} depois. Indices distintos andam juntos; mesmo indice espera. Vale a regra curta: pouco e fixo usa \texttt{critical(name)}; N variavel usa vetor de \texttt{locks} (\texttt{omp\_lock\_t}).
\subsection{Sobre o tempo: por que sobe com threads}
Cada insercao paga uma entrada em \texttt{critical}/\texttt{lock} mais a criacao de N \texttt{tasks}. Com N = 20000--50000 esse custo engole o trabalho util (um \texttt{malloc} + dois stores). Mais threads somam contencao e gerenciamento sem abrir paralelismo novo: as curvas sobem. Nome e \texttt{locks} removem o bloqueio cruzado. A fila na mesma lista fica — custo do dado compartilhado.
\section{Conclusao}
Foram geradas N \texttt{tasks} com \texttt{parallel+single}: \texttt{total==N} OK nas duas versoes. Foi sorteado com \texttt{rand\_r(\&seed)} privada: 24998/25002, $\approx$ N/2 por lista, sem a corrida do \texttt{rand()} global. Dois nomes isolam as listas na v1. \texttt{argv[2]} fixa N listas na v2, um \texttt{omp\_lock\_t} por lista. Nome nao escala porque sai resolvido em compilacao; \texttt{locks} indexaveis cobrem o N de execucao. Cada insercao paga uma entrada em \texttt{critical}/\texttt{lock}; privatizar nao cabe porque a lista e compartilhada.
\appendix
\newpage
\section{Codigo-fonte}
\lstinputlisting[caption={v1\_critical\_nomeado.c}]{v1_critical_nomeado.c}
\lstinputlisting[caption={v2\_locks.c}]{v2_locks.c}
\end{document}